\section{Especialización de la representación y decisión diagnóstica}
\label{sec:disc_representacion}
% Extensión estimada: 2 páginas, cinco párrafos (ficha 6.1).
% Fuentes: sec:res_vlm_controles (tab:res_vlm_controles), sec:res_vlm_dermlip (tab:res_dermlip, tab:res_dermlip_swin),
% tab:res_supervisados, sec:res_gen_sin_anclaje, sec:res_gen_ablaciones, sec:met_generacion; zeng2025mmskin,
% jiang2026multimodal, yan2026dermfmzero. Reemplaza las antiguas 6.1 (PI1) y 6.2.

Los \glspl{vlm} evaluados ocupan tres niveles en la clasificación de \emph{HAM10000}
(Sección~\ref{sec:res_comparacion}). Los generativos sin especialización en dermatología quedan cerca de 1/7 en sensibilidad macro, DermLIP sin ajuste
queda entre ellos y los supervisados, y DermLIP con una cabeza entrenada se acerca a los supervisados. La diferencia
entre esos niveles no es solo de grado, porque cada uno usa el modelo de una manera distinta. Esta sección discute qué
la explica: qué tiene DermLIP que los generativos no tienen, y qué les falta a estos para convertir lo que ven en un
diagnóstico.

Los controles de la Sección~\ref{sec:res_vlm_controles} permiten acotar las explicaciones posibles. Una de las hipótesis examinadas es que el bajo desempeño de los modelos generativos esté relacionado con la modalidad dermatoscópica, posiblemente menos representada durante su preentrenamiento. Para explorar esta posibilidad, se evalúa su comportamiento sobre fotografías clínicas de \emph{SkinCAP}, una modalidad más cercana a las imágenes de uso general. Las fotografías no mejoran su resultado, mientras que DermLIP distingue las clases en las dos modalidades. Además, con la imagen reemplazada por una gris, los generativos asignan todas las imágenes a una misma clase. La escasa variación de la sensibilidad macro al sustituir la imagen neutra por la imagen real sugiere que, bajo las instrucciones y estrategias evaluadas, la información visual tiene una influencia limitada sobre la decisión final del modelo. En conjunto, los resultados sugieren que las representaciones especializadas en dermatología podrían contribuir al mejor desempeño de clasificación. No obstante, las diferencias entre los modelos evaluados impiden establecer que esta sea la única causa de la ventaja observada.

El control de la representación indica dónde está esa especialización. Con las mismas tres cabezas, el
\emph{embedding} de DermLIP supera al de Swin preentrenado en \emph{ImageNet} por entre 0.15 y 0.35 de sensibilidad
macro (Tabla~\ref{tab:res_dermlip_swin}). Más llamativo es que esa representación, congelada y sin haber visto
\emph{HAM10000}, quede cerca de un clasificador ajustado de punta a punta sobre ese conjunto: con una \gls{svm}, su F1
macro es de 0.705, frente a 0.719 de Swin-Tiny, y con un \emph{linear probe} su sensibilidad macro supera a la de
Swin-Tiny, aunque con un F1 macro menor (Tablas~\ref{tab:res_dermlip} y~\ref{tab:res_supervisados}). Esas diferencias
no se compararon estadísticamente, de modo que solo indican que las dos representaciones quedan en un mismo rango. Una hipótesis compatible con estos resultados es que el entrenamiento contrastivo con pares de imágenes y descripciones dermatológicas favorezca representaciones que distingan características relevantes para la clasificación. Sin embargo, comprobar este mecanismo requeriría un análisis específico de la estructura del espacio de \emph{embeddings} y de su relación con los rasgos clínicos.

Mini-VLM, en cambio, extrae información de la imagen, pero, en estos experimentos, no la convierte en una clase. Después del
ajuste con \emph{SkinCaRe}, la descripción de Mini-VLM cambia con cada lesión, y sin la imagen su cobertura de
entidades clínicas cae a la mitad (Sección~\ref{sec:res_gen_ablaciones}). Sin un diagnóstico, sin embargo, el mismo
generador nombra la clase verdadera en menos de un reporte de cada cinco (Sección~\ref{sec:res_gen_sin_anclaje}). Es decir, la descripción del generador cambia con la imagen, mientras que, sin un diagnóstico, el reporte rara vez
nombra la clase verdadera. Una posible explicación es que las dificultades de clasificación no dependan exclusivamente de la representación visual, sino también de cómo el modelo generativo transforma esa información en una decisión entre categorías. Los resultados disponibles no permiten separar de manera concluyente ambas posibilidades.
Otra explicación, compatible con la anterior, es que \emph{SkinCaRe} agrupa las enfermedades en categorías distintas
de las siete de \emph{HAM10000} (Sección~\ref{sec:met_generacion}), de modo que el ajuste nunca le pidió al generador
esa correspondencia. Los experimentos no permiten separar las dos.

En el caso de SkinVL y con la configuración evaluada, la especialización dermatológica no fue suficiente. Ajustado con imágenes y textos
dermatológicos, sin un diagnóstico escribe casi siempre el mismo texto sobre una infección por hongos
(Sección~\ref{sec:res_gen_sin_anclaje}). Sus propios autores le miden a la variante que usa este trabajo, ajustada solo
con los libros de \emph{MM-Skin}, una exactitud del 3.5\,\% al clasificar \emph{HAM10000} sin ejemplos de ese conjunto, frente al
51.3\,\% de la variante que suma conjuntos de diagnóstico reformulados como preguntas~\cite{zeng2025mmskin}. La especialización de un generador parece quedar ligada, entonces, a las tareas con que se ajustó, y la de una representación
contrastiva se traslada a una tarea nueva con solo una cabeza. La literatura revisada es compatible con esta lectura: los generativos abiertos
obtienen exactitudes bajas también en consultas reales~\cite{jiang2026multimodal}, y por la vía contrastiva
DermFM-Zero supera a DermLIP en el protocolo de sus autores~\cite{yan2026dermfmzero}. Por lo tanto, la utilidad de un
\gls{vlm} para clasificar depende de cómo se lo use: como clasificador autónomo, ninguno de los evaluados alcanza a
los supervisados; como codificador de un modelo especializado, se acerca. El \emph{pipeline} aprovecha esa diferencia,
porque usa la representación de DermLIP para recuperar casos parecidos y toma el diagnóstico de un clasificador
supervisado. Queda por explicar por qué el generador, que no decide la clase, adopta sin más la que recibe.
